% Folder 146, batch 06 — archivists' pages 101–120.
% Pass: Opus 5.5 (claude-opus-5-5), 2026-10-02 — first pass, unchecked against the pages by a human.
%
% Pages 101 and 104 are blank and are not transcribed.
% Notation fixed for this batch: \mathbb{U} for the circle group (his
% hollow U), \mathbb{G}_m for his "Gm", \mathfrak{T} for his Gothic T
% (Teichmüller groupoids), \mathfrak{S} for the symmetric groups, G^\natural
% for the superscript he writes like a small "4" (maximal generalisation,
% p. 111), G^\flat for the superscript "b" of p. 102.
\documentclass[11pt,a4paper]{article}
\input{../preamble/grothendieck}

\folder{146}
\batch{6}
\pages{101}{120}
\foldertitle{Groupoïdes de Teichmüller et les S0 Tg,!ν, S0 TG : notes manuscrites (s.d.).}
\dating{[à partir de 1978]}
\watermark{Édition de démonstration}

\begin{document}

\page{102}
\note{la page s'ouvre au milieu d'un argument commencé avant ce lot : $SP_G$, $P_{G,C}$, $\widehat{M}_G$ et $S\mathfrak{T}_G$ y sont supposés connus.}

Le fibré $SP_G$, et tous les quotients $P_{G,C}$, \struck{\ill{}} sont des
fibrés sur $\widehat{M}_G$. \struck{Ils} \uncertain{Ils} sont déduisibles d'un
fibré encore plus fin, $SP_{G^{\flat}}$, de groupe $\mathbb{G}_m^{\vec{A}}$,
par l'hom. canonique $\mathbb{G}_m^{\vec{A}} \to \mathbb{G}_m^{A}$
(\uncertain{lié} sur $\widehat{M}_G$ ; où $SP^{\wedge}_{G^{\flat}}$). Donc
$S\mathfrak{T}_G$ déduisible de $S\mathfrak{T}_{G^{\flat}}$, de groupe
$\mathbb{Z}^{\vec{A}}$, par $\mathbb{Z}^{\vec{A}} \to \mathbb{Z}^{A}$.

\textit{NB} $\mathfrak{T}_G$ est une « \uncertain{cat.} » de groupoïdes
$R\mathfrak{T}_G$ par $\mathbb{Z}^{A}$ — elle est connue quand on connaît les
ext. correspondantes $\mathfrak{T}_{G,a}$ de $R\mathfrak{T}_G$
\add{$= R\mathfrak{T}_{G^{\flat}}$} \struck{par} $\mathbb{Z}$ — et il suffit
de connaître les ext. \add{$R\mathfrak{T}_{G,\vec{a}} = \mathfrak{T}_{G^{\flat},\vec{a}}$}
\struck{\ill{} de} $R\mathfrak{T}_G$ par $\mathbb{Z}$ pour $\vec{a} \in \vec{A}$.
\note{le mot entre guillemets, souligné, se lit « cat. » ; « ext. » irait mieux avec la suite.}

\marginal{NB \ill{} — une formule tracée en oblique dans la marge, où se lisent un $\mathfrak{T}$ et un produit $\prod$, reliée par une flèche à l'alinéa suivant.}

En définitive, on est réduit à construire les $S\mathfrak{T}_{g,I}$
($\simeq \mathfrak{T}^{!}_{g,\nu}$) \add{ext. de $\mathfrak{T}_{g,I}$ par $\mathbb{Z}^{I}$}
(ou encore les $S(\mathfrak{T}_{g,I})_i$ ($i \in I$) extension de
$\mathfrak{T}_{g,I}$ par $\mathbb{Z}$). De même, la construction des hom. de
généralisation $S\mathfrak{T}_G \to S\mathfrak{T}_{G'}$ se réduit au cas où
$G' = G^{\natural}$ \struck{\ill{}} ($G'$ \ill{} \uncertain{en cas où} \ill{}).

\page{103}
\note{en tête de page, un premier état du diagramme ci-dessous, barré de grands traits en zigzag ; il porte les mêmes sommets, avec $P_G (\simeq M_{G,G^{\natural}})$, $P_{G,C} (\simeq M_{G,G'})$, $RP_{G'} (= M_{G'})$, $RP_G (= M_G)$, un « $M_G \simeq$ » biffé devant ce dernier, et les mêmes mentions « cart ».}

\struck{NB les flèches horizontales}

$G \to G' \simeq G/C$ hom. de généralisation de graphes MD.

\begin{tikzcd}[column sep=small, row sep=small, nodes={font=\scriptsize}]
SP_G \arrow[r] \arrow[d, "\mathbb{G}_m^{I}"'] & SP_{G'} \arrow[r] \arrow[d, "\mathbb{G}_m^{I}"'] & SP_{G^{\natural}} \arrow[d, "\mathbb{G}_m^{I}"'] \\
P_G \, (\simeq M_{G,G^{\natural}}) \arrow[r] \arrow[d, "\mathbb{G}_m^{A'}"'] & P_{G'} \simeq M_{G',G^{\natural}} \arrow[r] \arrow[d, "\mathbb{G}_m^{A'}"'] & P_{G^{\natural}} = RP_{G^{\natural}} = M_{G^{\natural}} \\
P_{G,C} \simeq M_{G,G'} \arrow[r] \arrow[d, "\mathbb{G}_m^{C}"'] & RP_{G'} = M_{G'} & \\
RP_G = M_G & &
\end{tikzcd}

\note{les carrés sont marqués « cart » (cartésiens) sur la page.}

Les flèches horiz. ne sont définies que dans la catégorie homotopique. Elles
commutent aux opérations des groupes d'opérateurs $\mathbb{G}_m^{I}$ resp.
$\mathbb{G}_m^{A'}$ resp.
$\mathbb{G}_m^{\hat{A}'} = \mathbb{G}_m^{I} \times \mathbb{G}_m^{A'}$.

Par passage aux groupoïdes fondamentaux, on trouve :

\begin{tikzcd}[column sep=small, row sep=small]
S\mathfrak{T}_G \arrow[r] \arrow[d, "\mathbb{Z}^{I}"'] & S\mathfrak{T}_{G'} \arrow[r] \arrow[d, "\mathbb{Z}^{I}"'] & S\mathfrak{T}_{G^{\natural}} \arrow[d, "\mathbb{Z}^{I}"'] \\
\mathfrak{T}_G \arrow[r] \arrow[d, "\mathbb{Z}^{A'}"'] & \mathfrak{T}_{G'} \arrow[r] \arrow[d, "\mathbb{Z}^{A'}"'] & \mathfrak{T}_{G^{\natural}} \\
\mathfrak{T}_{G,C} \arrow[r] \arrow[d, "\mathbb{Z}^{C}"'] & R\mathfrak{T}_{G'} & \\
R\mathfrak{T}_G & &
\end{tikzcd}

\note{carrés marqués « cart » ; un signe biffé précède le $\mathfrak{T}_G$ de la deuxième ligne.}

Les flèches horizontales commutent aux op. de $\mathbb{Z}^{I}$ resp.
$\mathbb{Z}^{A'}$ resp. $\mathbb{Z}^{\hat{A}'} \simeq \mathbb{Z}^{I} \times \mathbb{Z}^{A'}$.

\note{sous un trait oblique qui sépare le bas de la page :}

\[
P_{G,C} = \prod_{a \in C \,/ M_G} P_{G,a}, \qquad
\mathfrak{T}_{G,C} = \Pi_1 P_{G,\ldots}
\]
\note{la seconde formule est coupée au bord de la feuille.}
\[
SP_G = P_{G,\hat{A}}, \quad RP_G = P_{G,\emptyset}, \quad SRP_G = P_{G,I}, \quad
P_G \simeq P_{G,A}.
\]
\note{le dernier indice de $SRP_G = P_{G,I}$ est à demi coupé au bord ; la lecture $I$ suit le carré ci-dessous.}

\begin{tikzcd}
SRP_G \arrow[d, "\mathbb{G}_m^{I}"'] & SP_G \arrow[l, "\mathbb{G}_m^{A}"'] \arrow[d, "\mathbb{G}_m^{I}"] \\
RP_G & P_G \arrow[l, "\mathbb{G}_m^{A}"]
\end{tikzcd}

\begin{tikzcd}
SR\mathfrak{T}_G \arrow[d, "\mathbb{Z}^{I}"'] & S\mathfrak{T}_G \arrow[l, "\mathbb{Z}^{A}"'] \arrow[d, "\mathbb{Z}^{I}"] \\
R\mathfrak{T}_G & \mathfrak{T}_G \arrow[l, "\mathbb{Z}^{A}"]
\end{tikzcd}

\note{chacun des deux carrés est marqué « cart. ».}

\page{105}
Soient $G \to G' \simeq G/C$ un hom. de généralisation de graphes MD, où
$C \subset A(G)$ (ens. des arêtes liées de $G$). Soit $X$ une courbe
algébrique sur $\mathbb{C}$, de type $G$. Je vais décrire une méthode pour
construire des courbes \add{$X'$} de type $G'$. On a
\[
C \subset A = \mathrm{Sing}(X), \qquad \text{soit } B = A \setminus C .
\]
Soit \struck{$\mathcal{C}$} $\widetilde{X}$ \add{de $\widetilde{B}$} le
normalisé de $X$ relativement à $C$, $\widetilde{C}$ (les images inverses
\add{de $B$} de $C$ dans $\widetilde{X}$). Choisissons pour tout
$x \in \widetilde{C}$ un disque fermé
$D_x \subset \widetilde{X} \setminus \widetilde{B}$, contenant $x$ dans son
intérieur, à bord \struck{tel que les} une courbe analytique réelle
(\uncertain{peut être choisi} — C$^1$ \struck{\ill{}} suffit), de façon que
les $D_x$ soient mutuellement disjoints. Il doit être vrai que l'ensemble des
choix possibles est \struck{tous} un espace contractile — i.e. que le choix est
« canonique » : la projection de la multiplicité \add{top.} formée des $X$ de
type $G$, \textit{munis} d'une famille $D_x$ (ou \ill{} des $X$
\uncertain{décomposés}) vers celle des $X$ de type $G$ sans plus, doit être une
équivalence d'homotopie. Soit
\[
X_1 = \widetilde{X} - \bigcup_{x \in \widetilde{C}} D_x^{\circ}
\]
qui est, \struck{en dehors de $X_1 \cap \mathrm{Sing}(\widetilde{X})$,} une
surface à bord holomorphe, avec
\[
\partial X_1 = \coprod_{x \in \widetilde{C}} \partial D_x .
\]
En tant que bord du disque holomorphe \struck{$D_x$}, de centre $x$,
$\partial D_x$ hérite une structure de $\mathbb{U}$-torseur

\page{106}
où on pose
\[
\mathbb{U} = \{ z \in \mathbb{C} \mid |z| = 1 \}
\]
(identifié au cercle de $\mathbb{U}$) — l'orientation de ce $\mathbb{U}$-torseur
\add{est opposée} de celle induite par $X_1$. \struck{\ill{}} La surface
holomorphe $X$ \struck{\ill{}} \add{avec la \ill{} qui} se récupère à partir de
$X_1$, surface holomorphe à bord $X_1$, et de la structure précédente
\add{$T_{\tilde{x}} = \partial_{\tilde{x}} X$} de $\mathbb{U}$-torseur sur
chacune des comp. connexes du bord $\partial X_1$, en notant que le cône sur
$\partial_{\tilde{x}} X$ \add{à bord} devient, grâce à cette structure, un
disque holomorphe, et \uncertain{par chirurgie holomorphe} on peut le recoller
à $X_1$, pour obtenir $X$, \ill{} les centres des disques
\uncertain{rapportés}.

Considérons, pour $x \in C$, les \add{deux} \struck{ensembles} \add{pts} $x'$
et $x''$ de $\widetilde{C}$ au-dessus de $x$, et \struck{l'espace} l'espace
\[
\mathrm{Isom}^{-}_{\mathbb{U}}(T_{x'}, T_{x''}) \simeq
\mathrm{Isom}_{\mathbb{U}}(T_{x'}^{-1}, T_{x''})
\]
des antiisomorphismes de $T_{x'}$ avec $T_{x''}$ (i.e. hom. de torseurs
compatibles avec l'\struck{isomorphisme} \add{antiautomorphisme}
$z \mapsto z^{-1} : \mathbb{U} \xrightarrow{\sim} \mathbb{U}$), qui
s'identifie, \ill{}, \uncertain{par symétrie}, au torseur
$T_{x'} \wedge T_{x''}$, ou \uncertain{encore} au $\mathbb{U}$-torseur
\[
T_x = \bigwedge_{\tilde{x} \in C_x} T_{\tilde{x}} .
\]

\marginal{NB $T_{\tilde{x}}$ s'identifie aussi \ill{} à $\widetilde{X}$ en
$\tilde{x}$, i.e. à $T^{*}_{\widetilde{X},\tilde{x}} / \mathbb{R}^{*+}$ (ce qui
est indép. du choix des $D_{\tilde{x}}$).}

Considérons enfin le produit
\[
P_C = P_{C,X} = \prod_{x \in C} T_x
\]
qui est un torseur sous $\mathbb{U}^{C}$. \struck{On} Pour
$\mathcal{X} = (X, (D_{\tilde{x}}))$

\page{107}
variable dans une multiplicité topologique convenable $M^C_G$, les
$P_{\mathcal{X}}$ forment un $\mathbb{U}^{C}_{M^C_G}$-torseur sur $M^C_G$, dit
$\widetilde{M}^C_G$, et on définit un \struck{application} morphisme de topos
\[
\widetilde{M}^C_G \longrightarrow M_{G'} .
\]
\note{sous $\widetilde{M}^C_G$, un « $M_{G'}$ » \uncertain{biffé}.}

\marginal{c'est l'espace d'un torseur évident \struck{\ill{}} $P_{G,C}$ ou $P_C$ sur \ill{}}

Revenons : $\mathcal{X} = (X, (D_{\tilde{x}}))$, on va définir, pour tout
$t = (t_x)_{x \in C} \in T_C$, une \struck{\ill{}} $X'$ de type $G'$, notée
$X'_t$. On l'obtient par chirurgie holomorphe : partir de $X_1$, en recollant
\add{$\forall x \in C$} les deux composantes $T_{x'}, T_{x''}$ du bord de $X_1$
correspondant à $x$, suivant l'antiisomorphisme $t_x$. Pour $\mathcal{X}, t$
variable, cela donne le morphisme de topos ci-dessus.

Soit $M_G$ \struck{\ill{}} la multiplicité holomorphe des courbes $X$ de type
$G$, on a
\[
M^C_G \longrightarrow M_G \quad \text{équiv. d'homotopie}
\]
donc

\begin{tikzcd}[column sep=small]
1 \arrow[r] & \pi_1(\mathbb{U}^{C}) \arrow[r] \arrow[d, no head, "\wr"'] & \pi_1(\widetilde{M}^C_G) \arrow[r] \arrow[dd] & \pi_1(M_G) \arrow[r] \arrow[d, no head, "\wr"] & 1 \\
 & \mathbb{Z}^{C} & & \pi_1(M^C_G) & \\
 & & \pi_1(M_{G'}) & &
\end{tikzcd}

Pour avoir $S\mathfrak{T}_G \to S\mathfrak{T}_{G'}$, on considère le fibré
$SP_{G'}$ sur $M_{G'}$, principal de groupe $\mathbb{U}^{\hat{A}'}$

\page{108}
(où $\hat{A}'$ est l'ens. des arêtes de $G'$) correspondant
\struck{au produit des}, pour une courbe $X'$ de type $G'$, à
\struck{\ill{}} \add{(contenant des \ill{} \struck{\ill{}} des $\hat{A}'$)}
\[
T_{X'} = \prod_{x \in \hat{A}'} T_{X',x}
\]
défini comme ci-dessus $T_{X,C}$, à la seule différence que pour les points
$x \in I'$, on prend $T_{X',x}$ = torseur tangent unitaire de $X'$ en $x$
$= T^{*}_{X',x} / \mathbb{R}^{*+}$. On a donc
\[
S\mathfrak{T}_{G'} = \pi_1(SP_{G'}) .
\]
D'autre part, soit \struck{\ill{} le fibré \ill{}} \add{$P_B$ le fibré}
\add{principal} canonique sur $M_G$, de groupe structural $\mathbb{U}^{B}$, de
sorte que l'on a
\[
SP_G = P_C \times_{M_G} P_B
\qquad \text{(isom. de $\mathbb{U}^{\hat{A}}$-torseurs sur $M_G$)}
\]
et considérons le diagramme ci-contre, cartésien : \uncertain{formé} \ill{}
déduit de la face de base (\uncertain{associée} à la formule précédente) par le
changement de base $M^C_G \to M_G$, qui est une équivalence d'homotopie,
donc \struck{toutes} les flèches verticales sont des équivalences
d'homotopie. On a donc

\begin{tikzcd}[column sep=small, row sep=small]
\widetilde{M}^C_G \arrow[rr] \arrow[dd] & & M^C_G \arrow[dd] & \\
 & \widetilde{P}^B_G \arrow[ul] \arrow[rr, dashed] \arrow[dd, dashed] & & P^B_G \arrow[ul, dashed] \arrow[dd, dashed] \\
P_C \arrow[rr] & & M_G & \\
 & SP_G \arrow[ul] \arrow[rr] & & P_B \arrow[ul]
\end{tikzcd}

\note{le sommet en haut à droite porte, au-dessus du $C$, une marque double qui pourrait être un second tilde ; la lecture $M^C_G$ suit la flèche verticale vers $M_G$.}

\page{109}
\[
S\mathfrak{T}_G \overset{\text{déf}}{=} \pi_1(SP_G) \xleftarrow{\ \sim\ } \pi_1(\widetilde{P}^B_G) .
\]
D'autre part, on a un diagramme cartésien évident

\begin{tikzcd}
\widetilde{P}^B_G \arrow[r] \arrow[d, "\text{$\mathbb{U}^B$-tors. sur $\widetilde{M}^C_G$}"'] & SP_{G'} \arrow[d] \\
\widetilde{M}^C_G \arrow[r] & M_{G'}
\end{tikzcd}

qui exprime simplement que pour $X'$ construit à l'aide de
$(X, (D_{\tilde{x}}))$, \struck{le} le $\mathbb{U}^{B}$-torseur $P_{B,X}$
\struck{\ill{}} est can. isom. au $\mathbb{U}^{B}$-torseur $P$ (on dirait même,
\textit{égal} ?) \ill{}. On trouve ainsi un morphisme

\begin{tikzcd}
\pi_1(\widetilde{P}^B_G) \arrow[r] \arrow[d, no head, "\wr"'] & \pi_1(SP_{G'}) \arrow[d, no head, "\wr"] \\
S\mathfrak{T}_G & S\mathfrak{T}_{G'}
\end{tikzcd}

Considérons le monoïde topologique multiplicatif $]0,1]$ (avec élément
unité $1$) et sa puissance $\vec{C}$-ième
\[
\mu_{\vec{C}} = \, ]0,1]^{\vec{C}}
\]
\add{il vaut mieux prendre $]0,1]^{C}$, sinon il y a des \ill{}}
\marginal{on prend $\rho_{\tilde{x}'} = \rho_{\tilde{x}''}$ pour $\tilde{x}', \tilde{x}''$ sur $x \in C$}
alors $\mu_{\vec{C}}$ opère sur la $M_G$-fibre $\widetilde{M}^C_G$ de
façons évidentes (homothéties sur les disques $D_{\tilde{x}}$),

\page{110}
donc aussi sur $\widetilde{P}^B_G$, \add{$\mathbb{U}^{B}$-tors.} image réciproque
de $P_B$ sur $M$\supplied{$_G$}. On a donc un \struck{hom.} morphisme de topos

\begin{tikzcd}[column sep=small]
\mu_{\vec{C}} \times \widetilde{M}^C_G \arrow[rr] \arrow[dr, dashed] & & M_{G'} \\
 & \widetilde{M}^C_G \arrow[ur, dashed, no head] &
\end{tikzcd}

s'insérant dans le diagramme cartésien

\begin{tikzcd}
\mu_{\vec{C}} \times P^B_G \arrow[r] \arrow[d, "\sim"'] & SP_{G'} \arrow[d] \\
\mu_{\vec{C}} \times M^C_G \arrow[r] & M_{G'}
\end{tikzcd}

L'\struck{hom.} morphisme $\mu_{\vec{C}} \times \widetilde{M}^C_G \to M_{G'}$
joue le rôle de l'\struck{inclusion} de $\mathcal{V}^{*}_{G,G'}$
\struck{dans} (\uncertain{l'ouvert} sur $M_{G'}$ du voisinage tubulaire de
$M_G$ dans $\widehat{M}_{G'}$) dans $M_{G'}$.

Les choses deviennent plus claires \struck{peut-être} (mais moins
intrinsèques) si on dispose d'une section de $M^C_G$ sur $M_G$ (y en a-t-il
toujours ??) — auquel cas on trouve un diagramme cartésien
\marginal{je pense que oui…}

\begin{tikzcd}
SP_G \times \mu_C \arrow[r] \arrow[d] & SP_{G'} \arrow[d] \\
P_C \times \mu_C \arrow[r] & M_{G'}
\end{tikzcd}

\page{111}
et ici, dans un sens très fort (qu'il vaudrait la peine peut-être
d'expliciter) $P_C \times \mu_C$ s'identifie à une forme de voisinage
tubulaire. \struck{et $P_C$ serait le bord}

Les cas les plus intéressants sont ceux où $G' = G^{\natural}$ (généralisation
maximale de $G$), ainsi que ceux où $G$ est de dim. modulaire $0$ ou $1$,
\struck{\ill{}} et où $G'$ est de dim. modulaire $\leq 2$ — notamment les
généralisations immédiates (correspondant au cas où $C$ est réduit à un
élément).

Cas où la dim. modulaire $d(G)$ de $G$ est nulle.
\marginal{i.e. les \uncertain{sommets} de $G$ sont de genre $0$, d'ordre $3$.}
Alors $X$ est défini à isom. unique près par la donnée de $G$, grâce au fait
qu'une droite projective munie d'un \struck{\ill{}} $I$-\uncertain{marquage},
où $\mathrm{card}\, I = 3$, est définie à isom. unique près. Donc la
multiplicité $M_G$ est la multiplicité ponctuelle. Considérons la droite
projective $\mathbb{P}^1_{\mathbb{C}}$ \struck{avec le groupe} avec les

\page{112}
trois pts marqués $0, 1, \infty$, avec le groupe d'automorphismes
holomorphes $\mathfrak{S}_3$, il y a une métrique \uncertain{positive}
\add{\ill{}} \struck{invariante} unique invariante par $\mathfrak{S}_3$. On va
prendre des \uncertain{systèmes} de disques $D_i$ ($i \in \{0,1,\infty\}$)
\textit{stables par} $\mathfrak{S}_3$, et \struck{fermés} limités par des
cercles de centre $i$ (\uncertain{pour} \struck{\ill{}} de la structure
riemannienne) — il revient au même de se donner un des $i$, disons
$i = \infty$, un cercle \add{de centre $i$} qui soit
\uncertain{invariant}* par la transposition \uncertain{qui échange} les deux
autres, et de plus \textit{petit} pour que ses \add{trois} transformés par
$\mathfrak{S}_3$ soient mutuellement disjoints.
\marginal{* c'est automatique}

\ill{} de stabilité, \struck{\ill{}} \uncertain{définissons} les disques
mutuellement disjoints. On prendra comme \add{\uncertain{modèle}} riemannien
une sphère \add{orientée} \struck{$\Sigma$} \struck{\ill{}} de rayon $1$,
$\Sigma = \Sigma_I$, un grand cercle fixé \add{$\Sigma_{\mathbb{R}}$}, et
\struck{\ill{}} \struck{la droite \ill{}} \struck{des cercles disjoints}
partie $I$ dessus formant un triangle équilatéral inscrit — le choix du
cercle dépend d'un paramètre $0 < \rho \leq 1$ (son rayon rectiligne). Pour
tout $i \in I$, le cercle $\Gamma_{i\rho}$ admet deux pts marqués sur
l'équateur, \struck{\ill{}} \add{une}

\drawing{dans la marge gauche, une sphère avec son équateur et un triangle inscrit en pointillé ; au-dessous, un cercle portant les points $\infty$ (en haut), $0$ et $1$, un petit cercle autour de $0$ et des segments en pointillé.}

\page{113}
bijection canonique
\[
\Gamma_{i\rho} \cap \Sigma_{\mathbb{R}} \simeq I \setminus \{i\}
\]
en associant à tout $x \in \Gamma_{i\rho} \cap \Sigma_0$ l'élément de
$I \setminus \{i\}$ le plus proche. Ces deux pts sont diamétralement opposés
sur $\Gamma_{i\rho}$, donc on trouve un isomorphisme canonique
\[
\Gamma_{i\rho} \simeq \text{\struck{$\Sigma$}}\, \mathbb{U} \wedge_{\{\pm 1\}} (I \setminus \{i\})
\]
i.e. $\Gamma_{i\rho}$ est déduit du $\{\pm 1\}$-torseur $I - \{i\}$ par
l'inclusion \struck{$(\Sigma \times \{\pm 1\})$} $\{\pm 1\} \hookrightarrow \mathbb{U}$.
\note{$\Sigma_{\mathbb{R}}$ et $\Sigma_0$ désignent ici, semble-t-il, le même grand cercle.}

\struck{La donnée d'} La courbe $X_G$ de type $G$ s'obtient en prenant la
normalisée
\[
\widetilde{X}_G = \coprod_{s \in G} \Sigma_{\hat{\vec{A}}_s}
\]
par recollement à partir de l'involution $\sigma$ sur
$\hat{\vec{A}} = \coprod_{s \in S} \hat{\vec{A}}_s$. On a pour tout
$x \in \hat{A}$, au-dessus de $s \in S$,
\[
T_x \simeq \mathbb{U} \wedge_{\{\pm 1\}} \{ \hat{\vec{A}}_s \setminus x \}
\quad \text{si } x \in I \simeq \hat{\vec{A}}^{\sigma}
\]
(en identifiant les éléments de $I$ à leurs images dans $\hat{A}$),
\[
T_x \simeq \mathbb{U} \wedge_{\{\pm 1\}} \bigwedge_{\tilde{x} \in \hat{\vec{A}}_x}
\bigl( \hat{\vec{A}}_{\sigma(\tilde{x})} \setminus \{\tilde{x}\} \bigr)
\quad \text{si } x \in A .
\]
\textit{NB} si $x$ est une boucle, on trouve $T_x \simeq \mathbb{U}$ iso. can.

\drawing{dans la marge gauche, un graphe : un sommet portant une boucle, un carré, des demi-arêtes terminées par des flèches, un triangle à gauche.}

\page{114}
On a donc
\[
P_G = \prod_{x \in \hat{A}} T_x
= \mathbb{U}^{\hat{A}} \wedge_{\{\pm 1\}^{\hat{A}}}
\overbrace{\prod_{a \in \hat{A}} \varepsilon_a}^{S\varepsilon_G}
\]
où pour $a \in \hat{A}$, $\varepsilon_a$ est l'ens. à deux éléments défini par
\[
\begin{cases}
\varepsilon_a = \hat{\vec{A}}_s - \{a\} & \text{si } a \in I \\
\varepsilon_a = \bigwedge_{\tilde{a} \text{ sur } a} \bigl( \hat{\vec{A}}_{\sigma(\tilde{a})} \setminus \tilde{a} \bigr) & \text{si } a \in A = \hat{A} \setminus I
\end{cases}
\]
et on trouve une \struck{\ill{}} \add{morphisme} canonique
\[
\mu^{\circ}_{\vec{A}} \times SP_G \longrightarrow SPM_{G^{\natural}}
\qquad \text{où } \mu^{\circ}_{\vec{A}} = (]0,1[)^{\vec{A}}
\]
définissant
\[
\pi_1(SP_G) = S\mathfrak{T}_G = \mathbb{Z}^{\hat{A}} \longrightarrow
\pi_1(SPM_{G^{\natural}}) = S\mathfrak{T}_{G^{\natural}}
\]
\note{« $SPM_{G^{\natural}}$ » est lu tel quel ; le $P$ et le $M$ sont serrés l'un contre l'autre.}
et, mieux, un hom. de groupoïdes fondamentaux. Mais
\struck{$\prod \varepsilon_a = \ldots \prod \varepsilon_G$} l'ensemble
$\varepsilon_G$ (de cardinal $2^{\mathrm{card}\, A}$) est contenu dans $SP_G$
\[
S\varepsilon_G \subset SP_G
\]
d'où
\[
\mu^{\circ}_{\vec{A}} \times S\varepsilon_G \longrightarrow PM_{G^{\natural}}
\]
et comme $\mu^{\circ}_{\vec{A}}$ est simplement connexe, on en conclut
\[
S\varepsilon_G \longrightarrow (S\mathfrak{T}_{G^{\natural}})
\]
(où $S\mathfrak{T}$ désigne le groupoïde fondamental de $SP_{G^{\natural}}$.)
\note{un signe biffé précède la parenthèse de $(S\mathfrak{T}_{G^{\natural}})$.}

\struck{Désignons par \ill{}} Désignons par $S_0\mathfrak{T}_G$ le
\add{s/}groupoïde de $S\mathfrak{T}_G = \Pi_1(\ldots)$ induit sur
$S\varepsilon_G$, on a donc
\[
S_0\mathfrak{T}_G \longrightarrow S\mathfrak{T}_{G^{\natural}} .
\]

\page{115}
Le groupoïde $S_0\mathfrak{T}_G$ est engendré par les objets de
$S\varepsilon_G = \prod_{a \in \hat{A}} \varepsilon_a$, et \struck{des
\ill{} morphismes \ill{}} des signes \add{$\varepsilon$}
\[
\alpha = (\alpha_a) \in \{\pm 1\}^{\hat{A}}
\]
et donc
\[
\xi \in \varepsilon_G
\]
une morphisme
\[
u_{\alpha,\xi} : \xi \longrightarrow \alpha.\xi = \xi'
\]
\add{défini par le chemin $u_{\alpha,\xi} : [0,1] \to P_G$}
\struck{\ill{}} \struck{$u_{\alpha,\xi}(t)_a =$} défini par
\[
u_{\alpha,\xi}(t)_a =
\begin{cases}
\xi_a = \xi'_a & \text{si } \alpha_a = +1 \\
(\exp(\pi i t))\, \xi_a & \text{si } \alpha_a = -1
\end{cases}
\]
avec les relations

\begin{tikzcd}
\xi \arrow[r, "u_{\alpha,\xi}"] \arrow[rr, bend right=30, "u_{\alpha\beta,\xi}\, n{(\alpha,\beta)}"'] & \alpha.\xi = \xi' \arrow[r, "u_{\beta,\xi'}"] & \beta\xi' = \beta\alpha\xi = \xi''
\end{tikzcd}

où
\[
n(\alpha,\beta) \in \mathbb{Z}^{\hat{A}}
\]
défini par \struck{$n(\alpha,\beta)_a = 0$ si \ill{}}
\[
n(\alpha,\beta)_a =
\begin{cases}
0 & \text{si } \alpha_a = 1 \text{ ou } \beta_a = 1 \\
1 & \text{sinon (i.e. } \alpha_a = \beta_a = -1)
\end{cases}
\]
On \struck{peut} a tenu compte, pour écrire ces relations, \struck{des} de la
$\mathbb{Z}^{\hat{A}}$-structure du groupoïde $S_0\mathfrak{T}_G$. Mais
indépendamment, pour ne pas l'utiliser, on introduit, pour $a \in \hat{A}$,
$\xi \in \varepsilon_G$
\[
u_{a,\xi} : \xi \longrightarrow \sigma_a \xi
\]
\marginal{NB $\sigma_a$ désigne l'involution de $S\varepsilon_G$ \uncertain{associée} à $a \in \hat{A}$}
et les relations suivantes

(1)

\begin{tikzcd}
\xi \arrow[r, "u_{a,\xi}"] \arrow[d, "u_{b,\xi}"'] & \sigma_a \xi \arrow[d, "u_{b,\sigma_a\xi}"] \\
\sigma_b \xi \arrow[r, "u_{a,\sigma_b\xi}"'] & \sigma_{ab}
\end{tikzcd}

commutatif ($\forall (\xi, a, b) \in S\varepsilon_G \times \hat{A} \times \hat{A}$ ;
il suffit de le prendre pour $a \neq b$)

\page{116}
(2) \struck{Les \ill{} $u_a$} soit $v^a_\xi$ \add{le composé}

\begin{tikzcd}
\xi \arrow[r, "u_{a,\xi}"] & a\xi \arrow[r, "u_{a,a\xi}"] & a^2 \xi = \xi
\end{tikzcd}

tous les $v^a_\xi$ commutent ($\xi$ fixé, $a \in \hat{A}$ variable).

\marginal{NB}
Le groupoïde libre sur \add{$\varepsilon_G$} \struck{les objets de} comme ens.
d'objets, et avec les générateurs $u_{a,\xi}$ et les relations (1), (2),
\struck{et} s'identifie bien à $S_0\mathfrak{T}_G$ (à isom. can. près).

Prenant $\mathfrak{T}_G$ au \struck{lieu} lieu de $S_0\mathfrak{T}_G$, on
trouve une construction \struck{\ill{}} \add{\uncertain{analogue}} en termes de
\struck{$\varepsilon_{G,A}$ \ill{} de $\varepsilon$}
\[
\varepsilon_G = \prod_{a \in A} \varepsilon_a \subset P_G = \prod_{a \in A} T_a
\]
qui est l'ens. des objets d'un groupoïde \add{\ill{}} ${}_0\mathfrak{T}_G$,
\struck{\ill{}} du groupoïde fondamental \struck{\ill{}}
$\mathbb{Z}^{A}$, \ill{}
\uncertain{comme} $S_0\mathfrak{T}_G$ à partir de $S\varepsilon_G$. On
\uncertain{en déduit}

\begin{tikzcd}
S_0\mathfrak{T}_G \arrow[r] \arrow[d] & S\mathfrak{T}_{G^{\natural}} \arrow[d] \\
{}_0\mathfrak{T}_G \arrow[r] & \mathfrak{T}_{G^{\natural}}
\end{tikzcd}

\note{les phrases autour de ${}_0\mathfrak{T}_G$ (« … du groupoïde fondamental $\mathfrak{T}_G$ … $\mathbb{Z}^{A}$ … ») sont en partie biffées et rapides ; seuls les termes mathématiques sont sûrs.}

\page{117}
\textit{Dimension modulaire 1}. Tous les sommets de $G$, sauf un seul $s_0$,
sont de genre $0$ et d'ordre $3$, tandis que $s_0$ est soit de genre $0$,
ordre $4$, soit de genre $1$, ordre $1$.

\drawing{deux esquisses du sommet $s_0$ : à gauche, quatre arêtes issues de $s_0$ ; à droite, une seule arête aboutissant à $s_0$, marquée $1$.}

Traitons d'abord le premier cas. Soit $J = \hat{\vec{A}}_{s_0}$, ens. de
cardinal $4$. On a \uncertain{un} iso canonique
\[
M_G \simeq M_{0,J}
\qquad \bigl( \simeq \mathbb{U}_{0,3} \simeq \mathbb{P}^1_{\mathbb{C}} \setminus \{0,1,\infty\} \bigr),
\]
\struck{\ill{}} isomorphisme qui dépend d'une \uncertain{numérotation} \ill{}.
\note{l'ensemble $\hat{\vec{A}}_{s_0}$ est noté d'une majuscule cursive, entre $I$ et $J$ ; on l'écrit $J$ comme à la p.~118.}

\struck{$SP_G \simeq$}

La donnée d'une courbe $X$ de type $G$ équivaut à celle d'une \struck{sphère}
droite projective $Y$ \add{$/J$-marquée}.

On a
\[
M_{0,J} = \text{\struck{$M_{0,4}$}}\; \mathrm{Rep}(J) \wedge_{\mathfrak{S}_4} M^{!}_{0,4}
= \mathrm{Rep}(J) \wedge_{\mathfrak{S}_4} \mathbb{U}_{0,3}
\]
\marginal{NB $\mathrm{Rep}(J) = \mathrm{Bij}(I_4, J)$, où $I_4 = \mathbb{N} \cap [1,4]$}
où $\mathfrak{S}_4$ opère sur $\mathbb{U}_{0,3}$ par l'intermédiaire de
l'homomorphisme
\[
\mathfrak{S}_4 \longrightarrow \mathfrak{S}_3
\]
et de l'action du groupe $\mathfrak{S}_3$ — ici $\mathfrak{S}_3$ est identifié
au \add{groupe} \struck{\ill{}} des permutations de l'ens. des partitions de
type $2,2$ de l'ens. à $4$ éléments $[1,4] \cap \mathbb{N} = I_4$, ces
partitions étant numérotées de $1$ à $3$, la partition $P_i$
($1 \leq i \leq 3$) \struck{\ill{}} étant \add{formée} de $\{i,4\}$ et de son
complémentaire \struck{\ill{}} $I_3 \setminus \{i\}$ dans $I_4$. On aura, si
$T_i$ ($i \in J$) sont les fibrés tangents unitaires de $Y$ aux pts de
$J \simeq \hat{\vec{A}}_{s_0}$, pour \struck{$T$} $a \in A$ :

\page{118}
\[
T_a = \mathbb{U} \wedge_{\{\pm 1\}} \varepsilon_a
\quad \text{pour} \quad
\varepsilon_a =
\begin{cases}
\hat{\vec{A}}_{\sigma(x)} \setminus x & \text{si } x \in I_s \\
\bigwedge_{\tilde{x} \in \hat{\vec{A}}_x} \bigl( \hat{\vec{A}}_{\sigma(\tilde{x})} \setminus \{\tilde{x}\} \bigr) & \text{si } x \in A
\end{cases}
\]
\marginal{si $s_0$ n'est pas un sommet de $a$}

\struck{Si $s_0$} Si $s_0$ est un sommet de $a$, on est dans l'un des
\add{trois} cas ci-dessous, suivant que $a$ \struck{est une arête \ill{}}
\add{\ill{}} ; est une boucle. \struck{\ill{}} Mais dans tous les cas on a

\drawing{trois esquisses numérotées : (I) une arête $a$ terminée par une flèche, issue d'un sommet portant d'autres arêtes ; (III) une arête $a$ joignant un sommet $s$ à $s_0$, avec la mention « $s \neq s_0$ » ; (II) une boucle $a$ en $s_0$.}

\[
T_a = \bigwedge_{\tilde{a} \in A} T_{\tilde{a}}
\]
avec
\begin{itemize}
  \item $T_i$ dans le cas (I), où $i \in J = \hat{\vec{A}}_{s_0}$ est l'élément qui correspond à l'unique $\vec{a}$ sur $a$ ;
  \item $\bigwedge_{i \in J \text{ au-dessus de } a} T_i$ dans le cas (II) ;
  \item $T_i \wedge_{\{\pm 1\}} \bigl( \hat{\vec{A}}_s - \{\vec{a}\} \bigr)$ dans le cas (III), où $\vec{a}$ désigne l'arc d'origine $s \neq s_0$ porté par $a$.
\end{itemize}

On trouve donc
\[
SP_X = \prod_{a \in \hat{A}} T_a
\]
et avec l'explicitation des $T_a$ par les formules précédentes, on voit que
chaque $T_i$ ($i \in J$) intervient dans le produit une fois et une seule,
mais sous \uncertain{forme tordue} si $i \in J = \hat{\vec{A}}_{s_0}$
correspond à une \ill{} qui est lié et pas une boucle.

Soit donc $\hat{A}_1 \subset \hat{A}$ l'ens. des arêtes de $G$ qui ne sont pas
\struck{directes} \uncertain{incidentes à} $s_0$ \uncertain{comme extrémité},
\struck{\ill{} qui sont liées et pas des boucles \ill{}} \struck{qui dans le
cas contraire sont \ill{}}

\page{119}
\struck{Posons, pour $a \in A_1$, $\varepsilon_a =$ ce qui est défini page
précédente, si $s_0$ pas \ill{} de $a$}

et $J_1$ l'ens. des éléments de $J = \hat{\vec{A}}_{s_0}$ qui sont
\struck{\ill{}} des arcs liés pas boucles, $J_2 = J \setminus J_1$. On aura
alors, posant
\[
\varepsilon_{\hat{A}_1} = \prod_{a \in \hat{A}_1} \varepsilon_a ,
\]
\[
SP_X = \underbrace{\bigl( \mathbb{U}^{A_1} \wedge_{\{\pm 1\}^{A_1}} \varepsilon_{A_1} \bigr)}_{T_{A_1}} \times T^{Y}_{J} ,
\]
où
\[
T^{Y}_{J} = \prod_{a \in \hat{A} \setminus \hat{A}_1} T_a
\]
(produit d'au moins deux facteurs, au plus quatre facteurs)
\marginal{$A \setminus A_1$ = ens. des arêtes incidentes à $s_0$}
où $T^{Y}_{J}$ est défini directement \add{\ill{}} en termes de $Y$ et des
ens. $\hat{\vec{A}}_s \setminus \{\vec{a}\}$ relatifs aux \struck{sommets}
\add{arcs $a \in \hat{\vec{A}}_{s_0} = J$} qui sont liés et pas des boucles,
$s$ désignant l'extrémité \struck{distincte de $s_0$ et adjacente à $s_0$}
\ill{} — \add{plus précisément} en termes de la \uncertain{façon} \ill{} comme
le groupage déduit des $J$ en \add{paires} (disparaissant dans le cas de deux
paires)), quand il y a — une ou deux — boucles en $G$ en $s_0$. De
\uncertain{façon} précise, par la structure de $G$ dans $J$ provient, on a
\[
J = J_{\mathrm{I}} \amalg J_{\mathrm{II}} \amalg J_{\mathrm{III}}
\]
où $J_{\mathrm{II}}$ est muni d'une involution sans pts fixes \add{$\sigma$}
et où $J_{\mathrm{III}}$ est muni d'un \uncertain{revêtement} de degré $2$
$\tilde{J}_{\mathrm{III}} \to J_{\mathrm{III}}$ \add{\uncertain{cette} structure
de $J$}. On a alors, \ill{},
\[
T^{Y}_{J} = T^{Y}_{J_{\mathrm{I}}} \times T^{Y}_{J_{\mathrm{II}}} \times T^{Y}_{J_{\mathrm{III}}}
\]

\page{120}
avec
\[
T^{Y}_{J_{\mathrm{I}}} = \prod_{i \in J_{\mathrm{I}}} T_i
\]
\[
T^{Y}_{J_{\mathrm{II}}} =
\underbrace{\prod_{\alpha \in J_{\mathrm{II}}/\sigma}}_{0,\,1 \text{ ou } 2 \text{ facteurs}}
\;
\underbrace{\bigwedge_{\substack{i \in J_{\mathrm{II}} \\ \text{sur } \alpha}} T_i}_{2 \text{ facteurs}}
\]
\[
T^{Y}_{J_{\mathrm{III}}} = \prod_{i \in J_{\mathrm{III}}} T_i \wedge_{\{\pm 1\}} \tilde{J}_{\mathrm{III},i}
\]
Il est clair que pour $Y$ variable, les $T^{Y}_{J_{\mathrm{I}}}$,
$T^{Y}_{J_{\mathrm{II}}}$, $T^{Y}_{J_{\mathrm{III}}}$ forment des
\add{$\mathbb{U}$-torseurs} \struck{fibrés} sur $M_{0,J}$, notés
$T_{\mathrm{I}}, T_{\mathrm{II}}, T_{\mathrm{III}}$, dont le produit fibré sera
noté $T_{A_2}$, où
\[
A_2 = A \setminus A_1 \simeq J_{\mathrm{I}} \amalg (J_{\mathrm{II}}/\sigma) \amalg J_{\mathrm{III}} .
\]
On aura donc
\[
SP_G \simeq
\underbrace{T_{\hat{A}_1}}_{\text{$\mathbb{U}^{\hat{A}_1}$-torseur}}
\times
\underbrace{T_{\hat{A}_2}}_{\text{$\mathbb{U}^{A_2}$-torseur sur $M_{0,J}$}}
\]

Nous allons maintenant décrire un \ill{} du groupoïde fondamental de $SP_G$,
i.e. de $S\mathfrak{T}_G$, par le choix de « pts bases » bien choisis. On va
s'intéresser aux \uncertain{$5$} pts de \uncertain{ceux} $Y$ qui correspondent
aux pts de $M_{0,J}$ qui sont des pts de ramification pour l'action du groupe
\uncertain{produit} $\mathfrak{S}_J$, de $\mathfrak{S}_{J'}$ \uncertain{où}
$J' \in \ldots$ ($3$ élts) est l'ens. des trois partitions de
\note{la phrase se poursuit au-delà de ce lot.}

\end{document}
